% Dinâmica: 2 blocos num plano inclinado + aceleração + forças
% Faro, qui 08 jan 2026 07:00:28 
% Written by Tordar 
% pstricks2pdf.sh plano_theta02
\documentclass{article}
\pagestyle{empty}

\usepackage{pst-all} 
\usepackage{pstricks-add}
\usepackage{pst-slpe}
\usepackage{graphicx}
\input{apnum}

% Bloco
\newcommand{\bloco}[5][]{% #1=options, #2=cx, #3=cy, #4=totalWidth, #5=height
  % Left frame (half of total width, full height)
  \psframe[linewidth=1pt,linecolor=gray,fillstyle=slope,slopebegin=white,slopeend=gray,slopesteps=15]
    (!#2 #4 2 div sub #3 #5 2 div sub)                     % lower-left of total area
    (!#2 #3 #5 2 div add)                                  % upper-right of left frame (center x, top)  
  % Right frame (half of total width, full height)
  \psframe[linewidth=1pt,linecolor=gray,fillstyle=slope,slopebegin=gray,slopeend=white,slopesteps=15]
    (!#2 #3 #5 2 div sub)                                  % lower-left of right frame (center x, bottom)
    (!#2 #4 2 div add #3 #5 2 div add)                     % upper-right of total area
  \psframe[#1](!#2 #4 2 div sub #3 #5 2 div sub)(!#2 #4 2 div add #3 #5 2 div add) 
}
% Roldana
\newcommand{\roldana}[5]{% x, y, radius, hole-radius, color
 \pscircle[linewidth=0.5mm,linecolor=lightgray,dimen=middle,fillstyle=radslope,slopebegin=gray,slopeend=white,slopesteps=50](#1,#2){#3}
  \ifnum\pdfstrcmp{#4}{0}=0\relax\else
    \pscircle[fillcolor=white,fillstyle=solid,linewidth=0.8pt](#1,#2){#4}
  \fi
  \pscircle[linewidth=1pt,linecolor=#5!50!black](#1,#2){#3}
}

\begin{document}

% Altura e largura do bloco 1: 
\evaldef\wbi{2}
\evaldef\hbi{1}
% Altura e largura do bloco 2: 
\evaldef\wbii{1.5}
\evaldef\hbii{1.5}
% Raio da roldana: 
\evaldef\ther{\hbi/2}
% Coordenadas do plano inclinado: 
\evaldef\xi{0}
\evaldef\xii{6}
\evaldef\yii{5}
% Inclinação do plano inclinado: 
\evaldef\thetheta{\ATAN{\yii/\xii}}%\apROUND\thetheta{6}
\evaldef\thethetad{\thetheta*180/\PI}\apROUND\thethetad{6}
% Coordenadas do bloco 1: 
\evaldef\xbi{\xii/2-\hbi/2*\SIN{\thetheta}}
\evaldef\ybi{\yii/2+\hbi/2*\COS{\thetheta}}
\evaldef\xpi{\xii/2}
\evaldef\ypis{\yii/2}
\evaldef\ypie{\ypis-2}
\evaldef\xvpi{\xpi+0.5}
\evaldef\yvpi{\ybi/2}
\evaldef\xfate{\wbi/2}
\evaldef\xfats{\xfate+1.5}
\evaldef\yfat{-\hbi/2}
% Coordenadas da aceleração: 
\evaldef\yabi{2}
\evaldef\xabis{-0.5}
\evaldef\xabie{ 0.5}
\evaldef\yvabi{\yabi-0.5}
% Coordenadas da roldana: 
\evaldef\xr{\xii+1}
\evaldef\yr{\xr*\TAN{\thetheta}}
% Coordenadas auxiliares: 
\evaldef\xfi{\xr-\hbi*\SIN{\thetheta}/2}
\evaldef\yfi{\yr+\hbi*\COS{\thetheta}/2}
\evaldef\ymi{0.85*\hbi}
\evaldef\ymr{\yr+1.75*\ther}
\evaldef\xtheta{2.5*\COS{\thetheta/2}}
\evaldef\ytheta{2.5*\SIN{\thetheta/2}}
% Coordenadas do bloco 2: 
\evaldef\xbii{\xr+\ther}
\evaldef\ybii{\yii/2}
\evaldef\xmii{\xbii+0.85*\wbii}
\evaldef\xabii{\xbii+2.5}
\evaldef\xvabii{\xbii+2}
\evaldef\yabiis{\ybii+0.5}
\evaldef\yabiie{\ybii-0.5}
\evaldef\ypii{\ybii-2}
\evaldef\xvpii{\xbii+0.75}
\evaldef\yvpii{\ybii-1.5}
\evaldef\ytii{\ybii+2}
\evaldef\yvtii{\ybii+1.5}

% Let's go:
\begin{pspicture}(0,0)(10,10)
% Chão: 
%\psframe[linewidth=0pt,linecolor=lightgray,dimen=inner,fillstyle=hlines*,fillcolor=lightgray,hatchangle=45](0,0)(10,2) 
% Plano inclinado: 
\pspolygon[linewidth=2pt](\xi,0)(\xii,\yii)(\xii,0)
% Arco do ângulo: 
\psarc[linewidth=3pt,arcsepA=3pt,arcsepB=3pt]{->}(0,0){2}{0}{\thethetad}
\psline[linewidth=2pt,linecolor=black](\xbi,\ybi)(\xfi,\yfi) % fio bloco 1 - periferia da roldana
\psline[linewidth=2pt,linecolor=black](\xbii,\ybii)(\xbii,\yr) % fio bloco 2 - periferia da roldana
% Bloco 1: 
\rput{\thethetad}(\xbi,\ybi){
\rput{0}(\xi,\ymi){\scalebox{1.5}{$m_1$}}
\rput{0}(\xi,\yabi){\scalebox{1.5}{$\mathbf{a}$}}
\rput{0}(2,0.5){\scalebox{1.5}{$\mathbf{T}_1$}}
\rput{0}(2,-1){\scalebox{1.5}{$\mathbf{F}_{at}$}}
\psline[linewidth=3pt,linecolor=blue]{->}(\xabis,\yvabi)(\xabie,\yvabi)
\psline[linewidth=3pt,linecolor=blue]{->}(\xfats,\yfat)(\xfate,\yfat) % Força de atrito
\psline[linewidth=3pt,linecolor=blue]{->}(\xi,0)(2.5,0)
\bloco[linewidth=2pt,linecolor=black]{\xi}{0}{\wbi}{\hbi}
}
\psline[linewidth=3pt,linecolor=blue]{->}(\xpi,\ypis)(\xpi,\ypie)
% Bloco 2:
\psline[linewidth=3pt,linecolor=blue]{->}(\xbii,\ybii)(\xbii,\ypii) % Peso 
\psline[linewidth=3pt,linecolor=blue]{->}(\xbii,\ybii)(\xbii,\ytii) % Tensão
\bloco[linewidth=2pt,linecolor=black]{\xbii}{\ybii}{\wbii}{\hbii}
\psline[linewidth=3pt,linecolor=blue]{->}(\xvabii,\yabiis)(\xvabii,\yabiie)
% Roldana: 
\roldana{\xr}{\yr}{\ther}{0.05}{gray} 
\psline[linewidth=2pt,linecolor=black](\xii,\yii)(\xr,\yr) % fio plano - centro roldana
% Texto: 
\rput{0}(\xmii,\ybii){\scalebox{1.5}{$m_2$}}
\rput{0}(\xr,\ymr){\scalebox{1.5}{$m_R$}}
\rput{0}(\xtheta,\ytheta){\scalebox{1.5}{$\theta$}}
\rput{0}(\xabii,\ybii){\scalebox{1.5}{$\mathbf{a}$}}
\rput{0}(\xvpi,\yvpi){\scalebox{1.5}{$\mathbf{P}_1$}}
\rput{0}(\xvpii,\yvpii){\scalebox{1.5}{$\mathbf{P}_2$}}
\rput{0}(\xvpii,\yvtii){\scalebox{1.5}{$\mathbf{T}_2$}}
\end{pspicture}

\end{document}
